\subsection{A structural twelve-arrow exclusion with arrows in both directions}

\begin{lemma}[The minimal ideal of a finite strongly connected category]
\label{C:finite-category-Rees}
Let $C$ be a nonempty finite category in which every hom-set is
nonempty. There are a finite group $G$, positive integers $r_i,c_i$
for each object $i$, and a least nonempty two-sided ideal $K$ of
arrows, with
\[
 K(i,j)\cong R_j\times G\times C_i,
 \qquad |R_j|=r_j,\quad |C_i|=c_i.
\]
After choosing these coordinates, composition has the form
\[
 (r_k,g,c_j)(r_j,h,c_i)
       =(r_k,g\,p_{j,c_j,r_j}\,h,c_i),\qquad p_{j,c_j,r_j}\in G.
\]
In particular,
\[
 |K|=|G|\Bigl(\sum_i r_i\Bigr)\Bigl(\sum_i c_i\Bigr).
\]
If every $r_i=1$, or every $c_i=1$, then $BC\simeq BG$.
If $1_i\in K(i,i)$, object $i$ is a retract of every object of $C$.
\end{lemma}
\begin{proof}
Fix an object $a$, choose the least nonempty ideal of the finite
monoid $\operatorname{End}(a)$, and choose an idempotent $e$ in
that ideal. Such an idempotent is a sufficiently large power of
any element of the ideal. The monoid
$G=e\operatorname{End}(a)e$ is a group. Indeed, for every
$z\in G$, minimality gives $e=u z v$ for some endomorphisms
$u,v$; after sandwiching, this is
$e=(eue)z(eve)$ in the finite monoid with identity $e$.
In a finite monoid, every factor of a product equal to the
identity is a unit, so $z$ is a unit.

Put
\[
 U_j=C(a,j)e,\qquad V_i=eC(i,a),\qquad K(i,j)=U_jV_i.
\]
These sets are nonempty and the $K(i,j)$ form an ideal.
Every nonempty ideal of $C$ meets $\operatorname{End}(a)$:
compose any one of its arrows with arrows from and to $a$.
Its intersection with that endomorphism monoid consequently
contains $e$, and it contains every $U_jV_i$. Thus $K$ is the
least nonempty ideal.

The right $G$-action on $U_j$ and the left $G$-action on $V_i$
are free. For example, if $u\in U_j$ and $v\in V_j$, then
$vu\in G$, so $ug=u$ implies $(vu)g=vu$ and hence $g=e$.
The other action is analogous. Moreover, multiplication gives
a bijection
\[
 U_j\times_G V_i\longrightarrow K(i,j).
\]
Surjectivity is the definition. If $uv=u'v'$, choose $w\in U_i$
and adjust it by an element of $G$ so that $vw=e$.
Then $u=u'(v'w)$, with $v'w\in G$. Applying a left inverse
of $u'$ to the original equality gives the corresponding
relation between $v$ and $v'$. This is precisely the balanced
$G$-equivalence relation. Choosing orbit representatives now
gives the stated coordinates; the intervening products
$v_{j,c}u_{j,r}\in G$ give the displayed sandwich entries.

Adjoin the retract object represented by $e$, inside the
idempotent completion, and call the resulting category $C^+$.
This does not change the classifying space. To recall the
elementary Quillen-A argument, write
$u:e\to a$, $v:a\to e$ with $vu=1_e$.
In $(e\downarrow C)$, the object $u$ admits a natural
transformation from its constant functor to the identity,
whose component at $f:e\to i$ is $fv:a\to i$.
This contracts the comma category; commas for objects already
in $C$ have initial identity objects.

The one-object full subcategory at $e$ has endomorphism group
$G$. If every $r_i=1$, then every $U_i=C^+(e,i)$ is a free
transitive right $G$-set. The comma category $(G\downarrow i)$
is its translation groupoid and is contractible. Quillen's
Theorem A gives $BG\simeq BC^+\simeq BC$.
If every $c_i=1$, apply the opposite argument.

Finally, if $1_i=uv\in K(i,i)$, then $vu\in G$ is idempotent
and therefore equals $e$. For any object $j$, choose
$u_j\in U_j$, $v_j\in V_j$ and adjust so that $v_ju_j=e$.
The arrows $u_jv:i\to j$ and $uv_j:j\to i$ have composite
$uv_ju_jv=uv=1_i$, proving the retract assertion.
\end{proof}

\begin{lemma}[The rectangular minimal ideal Tor comparison for a category]
\label{C:category-band-Tor}
Suppose the group in Lemma~\ref{C:finite-category-Rees} is
trivial. Put $R=\mathbb QC$ and $A=R/\mathbb QK$, and choose
the minimal-corner idempotent $e$ used in that proof.
Let $\mathbf 1_\ell$ and $\mathbf 1_r$ be the constant left
and right $R$-modules, with one copy of $\mathbb Q$ at every
object and every arrow acting as the identity between its
source and target copies. Define
\[
 V=\ker(Re\longrightarrow\mathbf1_\ell),\qquad
 W=\ker(eR\longrightarrow\mathbf1_r),
\]
where the maps send each arrow basis element to one in its
object component. Then $V,W$ are $A$-modules and, for $q\geq2$,
\[
 H_q(BC;\mathbb Q)
       \cong\operatorname{Tor}_{q-2}^{A}(W,V).
\]
Here $\dim V=\sum_i(r_i-1)$ and
$\dim W=\sum_i(c_i-1)$.
\end{lemma}
\begin{proof}
The module $Re$ has $r_i$ basis vectors at object $i$, and
$eR$ has $c_i$ basis vectors there. Their augmentations are
therefore surjective and have the stated kernel dimensions.
Every arrow in $K$ acts constantly on the corresponding row
or column set, so $K$ annihilates both augmentation kernels.
The coordinate description also gives an $R$-bimodule
isomorphism
\[
 Re\otimes_{\mathbb Q}eR\xrightarrow{\ \mathrm{multiplication}\ }
 \mathbb QK.
\]
In fact its arrow-basis tensors map bijectively to the row-column
basis of $K$. It follows that $\mathbb QK$ is projective on
both sides.

The normalized category bar computes
$H_q(BC;\mathbb Q)=\operatorname{Tor}_q^R(\mathbf1_r,\mathbf1_\ell)$.
Dimension shifting first through $Re$ and then through $eR$
gives
\[
 H_q(BC;\mathbb Q)\cong\operatorname{Tor}_{q-2}^R(W,V)
 \qquad(q\geq2).
\]
At $q=2$, the second shift uses the exact sequence
\[
 0\longrightarrow\operatorname{Tor}_1^R(\mathbf1_r,V)
 \longrightarrow W\otimes_R V\longrightarrow eR\otimes_R V.
\]
The last term is $eV=0$, so this is indeed the same formula
with ordinary tensor product in degree zero.

There is no extra derived change-of-rings term here.
The projective left resolution
$0\to\mathbb QK\to R\to A\to0$ gives
\[
 W\otimes_R\mathbb QK
 \cong We\otimes_{\mathbb Q}eR=0,
\]
and therefore $\operatorname{Tor}_j^R(W,A)=0$ for $j>0$,
while $W\otimes_R A=W$. The change-of-rings spectral sequence
\[
 \operatorname{Tor}_i^A\bigl(\operatorname{Tor}_j^R(W,A),V\bigr)
 \Longrightarrow\operatorname{Tor}_{i+j}^R(W,V)
\]
has only its $j=0$ row. This proves the asserted comparison.
\end{proof}

\begin{proposition}[The nontriangular two-object case through twelve arrows]
\label{C:twelve-strongly-connected}
No category with two objects, nonempty hom-sets in both
directions, and at most twelve arrows has classifying space
homotopy equivalent to $\mathbb CP^2$.
\end{proposition}
\begin{proof}
Suppose otherwise, and use the notation of
Lemma~\ref{C:finite-category-Rees}. If an object identity
belongs to $K$, that object is a retract of the other object.
Deleting it preserves the classifying space by the same comma
contraction used in that lemma, and leaves at most eleven
arrows. This contradicts Proposition~\ref{C:CP2-lower12}.
Thus both identities lie outside $K$.

If all row numbers, or all column numbers, are one, the same
lemma gives $BC\simeq BG$. A finite group has no positive
rational homology, so this cannot be $\mathbb CP^2$.
Otherwise both total row and column numbers are at least
three. The arrow count now gives
\[
 |G|\,(r_a+r_b)(c_a+c_b)+2\leq12.
\]
It forces $G=1$ and
\[
 r_a+r_b=c_a+c_b=3.
\]
The ideal therefore has exactly nine arrows, and there is at
most one additional arrow beyond its arrows and the two
identities. Up to exchanging the objects and taking opposites,
there are just two configurations.

First suppose the doubled row and column occur at the same
object $a$, so $(r_a,r_b)=(c_a,c_b)=(2,1)$.
The minimal hom-cardinalities are
\[
 \begin{pmatrix}4&2\\2&1\end{pmatrix}.
\]
At least one cross hom-set has exactly its two minimal arrows.
Put $M=\operatorname{End}(a)$; its cardinality is at most six.
If $C(b,a)=K(b,a)$, this left $M$-set is isomorphic to $Mf$
for any idempotent $f\in K(a,a)$: both consist of the two
row coordinates with a fixed column. The comma category
$(b\downarrow M)$ is its action category and is contractible.
For clarity, the action category of $M$ on $Mf$ contracts
by the natural transformation from the constant object $f$
to the identity, whose component at $z\in Mf$ is the arrow
labelled $z$, since $zf=z$.
The other comma category has an initial identity object.
Quillen's Theorem A gives $BM\simeq BC$.
If instead $C(a,b)=K(a,b)$, use the right $M$-set $fM$ and
the opposite comma argument. In either case a monoid of at
most six elements would realize $\mathbb CP^2$, again
contradicting Proposition~\ref{C:CP2-lower12}.

Second suppose the doubled row and column occur at different
objects. Choose names so that $(r_a,r_b)=(2,1)$ and
$(c_a,c_b)=(1,2)$. The minimal endomorphism cardinalities
are two at each object, and the cross cardinalities are
$|K(a,b)|=1$, $|K(b,a)|=4$.
If no additional arrow is placed in $C(a,b)$, the singleton
return criterion, Lemma~\ref{C:singleton-return}, makes $BC$
contractible. The only remaining possibility thus has exactly
twelve arrows: the nine arrows of $K$, two identities, and
one extra arrow $t:a\to b$.

In this last case $A=\mathbb QC/\mathbb QK$ has basis
$1_a,1_b,t$ and is precisely the path algebra of one arrow
$a\to b$. It is hereditary. Explicitly, for a representation
$f:X_a\to X_b$, its projective resolution is
\[
 0\longrightarrow P_b\otimes X_a
 \longrightarrow(P_a\otimes X_a)\oplus(P_b\otimes X_b)
 \longrightarrow X\longrightarrow0,
\]
with the maps at vertex $b$ given by
$x\mapsto(x,-f(x))$ and $(x,y)\mapsto f(x)+y$.
Lemma~\ref{C:category-band-Tor} consequently gives
\[
 H_4(BC;\mathbb Q)=\operatorname{Tor}_2^A(W,V)=0.
\]
This is the final contradiction. The proof covers every
composition law satisfying the indicated minimal-ideal
structure; no twelve-arrow composition-table census is used.
\end{proof}
